\section{Theoretical Setup}
\label{sec:setup}

\subsection{Problem setting}

Fix a tie label $\tau\in\{-1,+1\}$, equal to the convention used for
$\operatorname{sign}(0)$ in the learning protocol, and define
\[
\operatorname{sgn}_{\tau}(z):=
\begin{cases}
+1,&z>0,\\
-1,&z<0,\\
\tau,&z=0.
\end{cases}
\]
Let $\mathcal X=\{-1,+1\}^n$, let $\Delta(\mathcal X)$ denote the set of
probability distributions on $\mathcal X$, and let
$\mathcal H\subseteq\{-1,+1\}^{\mathcal X}$.  Since $\mathcal X$ is
finite, $\mathcal H$ is finite.  Its deterministic dimension complexity is
\[
\operatorname{dc}(\mathcal H)
:=\min\left\{q\in\mathbb Z_{\geq0}:
\begin{array}{l}
\exists\Phi:\mathcal X\to\mathbb R^q\ \forall h\in\mathcal H\
\exists u_h\in\mathbb R^q\ \forall x\in\mathcal X,\\[1mm]
\operatorname{sgn}_{\tau}(\langle u_h,\Phi(x)\rangle)=h(x)
\end{array}
\right\}.
\]
We set $\operatorname{dc}(\varnothing)=0$ and use
$\mathbb R^0=\{0\}$ and $\langle0,0\rangle=0$.

Fix a fully connected, bias-free ReLU architecture with $n_0=n$, $n_L=1$,
positive integer widths $n_1,\ldots,n_{L-1}$, parameter matrices
$\theta_i\in\mathbb R^{n_i\times n_{i-1}}$, and parameter count
\[
S:=\sum_{i=1}^{L}n_i n_{i-1}.
\]
For $z_0=x$, set $z_i=\sigma(\theta_i z_{i-1})$ for $1\leq i<L$, where
$\sigma(a)=\max\{0,a\}$ coordinatewise, and set
$f_\theta(x)=\theta_Lz_{L-1}$.  Initialize independently according to
\[
(\theta_i^{(0)})_{ab}\sim\mathcal N(0,1/n_{i-1}).
\]
At ReLU kinks, $\nabla^{\mathrm{src}}$ denotes the fixed gradient or
subgradient convention of the protocol; it makes the recursion single-valued
without assuming that kinks are avoided.  Given $\mathcal D\in\Delta(\mathcal
X)$ and $h^\star\in\mathcal H$, draw
$x^{(0)},\ldots,x^{(T-1)}\stackrel{\mathrm{iid}}{\sim}\mathcal D$ and perform
\[
\theta^{(t+1)}=\theta^{(t)}-
\eta\nabla^{\mathrm{src}}_\theta
\ell\!\left(h^\star(x^{(t)})f_{\theta^{(t)}}(x^{(t)})\right),
\qquad
\ell(z)=\log(1+e^{-z}),
\quad 0\leq t<T.
\]
The latter-half aggregate and returned classifier are
\[
A_{\mathcal D,h^\star}(x)
:=\sum_{t=\lceil T/2\rceil}^{T}f_{\theta^{(t)}}(x),
\qquad
\widehat h_{\mathcal D,h^\star}(x)
:=\operatorname{sgn}_{\tau}(A_{\mathcal D,h^\star}(x)).
\]
For any binary classifier $g$, write
\[
\mathcal L_{\mathcal D,h}(g)
:=\Pr_{x\sim\mathcal D}[g(x)h(x)<0]
=\Pr_{x\sim\mathcal D}[g(x)\neq h(x)].
\]

\subsection{Assumptions}

\begin{assumption}[Source parameter regime]
\label{assump:source-regime}
The integers $n,L,T$ are positive, all network widths are positive integers
with $n_0=n$ and $n_L=1$, $\eta>0$, $0\leq\varepsilon<1/4$, and the
confident-map dimension $d\in\mathbb Z_{\geq0}$ is an exposed parameter.
\end{assumption}

\begin{assumption}[Universal exact-SGD success]
\label{assump:universal-sgd-success}
The architecture, $\eta$, and $T$ are fixed before the distribution and
target are chosen, and
\[
\forall\mathcal D\in\Delta(\mathcal X)\ \forall h^\star\in\mathcal H,
\qquad
\mathbb E_{\theta^{(0)},\,
x^{(0:T-1)}\stackrel{\mathrm{iid}}{\sim}\mathcal D}
\left[
\mathcal L_{\mathcal D,h^\star}
(\widehat h_{\mathcal D,h^\star})
\right]
\leq\varepsilon.
\]
The expectation is over exactly the Gaussian initialization and the one-sample
SGD draws described above; the loss, update, architecture, and returned
predictor are unchanged.
\end{assumption}

\begin{assumption}[Target-independent tie-resolved confident map]
\label{assump:tie-resolved-confident-map}
There exists one probability law $\mathcal P$ on maps
$\phi:\mathcal X\to\mathbb R^d$, selected before $\mathcal D$ and $h$, such
that
\[
\exists\mathcal P\ \forall\mathcal D\in\Delta(\mathcal X)\
\forall h\in\mathcal H,
\qquad
\Pr_{\phi\sim\mathcal P}\!\left[
\exists w\in\mathbb R^d\ \forall x\in\mathcal X,
\ \operatorname{sgn}_{\tau}(\langle w,\phi(x)\rangle)=h(x)
\right]\geq\frac12.
\]
Thus $\mathcal P$ is independent of the subsequently quantified distribution
and target, whereas $w$ may depend on $(\phi,h)$.  The event is exact and
tie-resolved; no nonzero margin or avoidance of zero scores is assumed.
\end{assumption}

\subsection{Formalized goal}

The objective is the following conditional deterministic amplification
statement with the explicit universal constant $7$:
\[
\operatorname{dc}(\mathcal H)\leq 7TSd.
\]
The constant must have no dependence on
$n,\mathcal H,L,(n_i)_{i=0}^L,S,\eta,T,\varepsilon,d$, or $\mathcal P$.
The proof must derive $\operatorname{VC}(\mathcal H)<2T$, the finite-class
count, and one deterministic common feature map.  If a separate result also
establishes an explicit polynomial bound $d\leq p(S,T)$ with no hidden
$n$- or $\eta$-dependence, the same theorem specializes to
$\operatorname{dc}(\mathcal H)\leq7TSp(S,T)$.
Assumption~\ref{assump:tie-resolved-confident-map}, and in particular such a polynomial
bound on $d$, is an additional conditional premise; it is not claimed to
follow from Assumption~\ref{assump:universal-sgd-success}.
